\documentclass[11pt]{report}

\usepackage[a4paper,margin=1in]{geometry}
\usepackage[T1]{fontenc}
\usepackage[utf8]{inputenc}
\usepackage{lmodern}
\usepackage{amsmath}
\usepackage{amssymb}
\usepackage{graphicx}
\usepackage[hidelinks]{hyperref}
\usepackage{titlesec}
\usepackage{fancyhdr}
\usepackage{xcolor}
\usepackage{tcolorbox}
\usepackage{setspace}
\usepackage{enumitem}
\usepackage{parskip}

\onehalfspacing

\definecolor{accent}{RGB}{40,80,160}
\definecolor{soft}{RGB}{240,245,255}

\titleformat{\chapter}
{\Huge\bfseries\color{accent}}
{\thechapter}{20pt}{}

\titleformat{\section}
{\Large\bfseries}
{\thesection}{12pt}{}

\titleformat{\subsection}
{\large\bfseries}
{\thesubsection}{10pt}{}

\pagestyle{fancy}
\fancyhf{}
\fancyhead[L]{Portfolio Analytics Handbook}
\fancyhead[R]{\thepage}

\newtcolorbox{callout}{
colback=soft,
colframe=accent,
boxrule=0.6pt,
arc=4pt
}

\title{
\Huge Portfolio Analytics Handbook\\
\vspace{6pt}
\Large Monte Carlo Portfolio Risk Simulator
}

\author{Portfolio Analytics Project}
\date{\today}

\begin{document}

\maketitle
\clearpage

\tableofcontents
\clearpage

\chapter{Foundations}

\section{Quick Start: How to Use This Simulator}
\begin{callout}
\textbf{Level.} Beginner
\end{callout}

This simulator helps you explore how a portfolio might behave under many possible future scenarios.

\textbf{A practical workflow.}

\begin{enumerate}
\item Choose the assets you want in the portfolio.
\item Set portfolio weights.
\item Choose simulation settings such as time horizon and number of paths.
\item Review the outputs.
\end{enumerate}

\begin{itemize}
\item simulated portfolio paths
\item return distribution
\item terminal value distribution
\item risk metrics
\item optimization results
\item historical backtesting
\item factor analysis
\item stress testing
\item investment projections
\end{itemize}

This app does not predict one exact future.

Instead, it builds many possible futures under a statistical model so you can study ranges, risks, and trade-offs.

\section{What This App Does}
\begin{callout}
\textbf{Level.} Beginner
\end{callout}

\textbf{This dashboard combines several portfolio analytics tools in one place.}

\begin{itemize}
\item Monte Carlo simulation for future portfolio scenarios
\item risk analysis using volatility, VaR, and CVaR
\item portfolio optimization and efficient frontier analysis
\item backtesting against historical data
\item benchmark comparison
\item CAPM-style market factor analysis
\item stress testing under adverse assumptions
\item long-term investment outcome projection
\end{itemize}

The goal is not only to compute numbers, but to help users understand what those numbers mean.

\section{Glossary}
\begin{callout}
\textbf{Level.} Beginner
\end{callout}

\begin{description}[leftmargin=3.4cm, style=nextline, font=\normalfont\bfseries]
\item[Portfolio] A collection of assets held together as one investment.
\item[Weight] The proportion of the portfolio allocated to a particular asset.
\item[Return] The percentage gain or loss over a period.
\item[Expected Return] The model-based average return across scenarios.
\item[Volatility] A measure of how much returns fluctuate.
\item[Covariance] A measure of how two assets move together.
\item[Correlation] A standardized version of co-movement between assets.
\item[Monte Carlo Simulation] A method that generates many random future scenarios.
\item[Benchmark] A reference portfolio or market index used for comparison, such as SPY.
\item[VaR (Value at Risk)] A lower-tail threshold for losses.
\item[CVaR (Conditional Value at Risk)] The average loss in the worst tail beyond VaR.
\item[Sharpe Ratio] A measure of return relative to risk.
\item[Beta] Sensitivity of a portfolio to market movements.
\item[Alpha] Return not explained by market exposure alone.
\item[$R^2$] The fraction of portfolio return variation explained by the factor model.
\item[Terminal Value] The ending portfolio value at the end of the simulation horizon.
\item[Efficient Frontier] A set of portfolios representing attractive return-risk combinations under the model.
\end{description}

\section{Portfolio Basics}
\begin{callout}
\textbf{Level.} Beginner
\end{callout}

A portfolio is a weighted combination of assets such as stocks or ETFs.

\begin{callout}
\textbf{Example.}

\begin{itemize}
\item 50\% SPY
\item 30\% QQQ
\item 20\% NVDA
\end{itemize}
\end{callout}

\textbf{Portfolio behavior depends on.}

\begin{itemize}
\item the expected return of each asset
\item the volatility of each asset
\item the correlation structure across assets
\item the weights chosen by the investor
\end{itemize}

A strong portfolio is not simply a list of good assets.

It is a design problem: how assets are combined matters just as much as which assets are selected.

\section{Risk and Return}
\begin{callout}
\textbf{Level.} Beginner
\end{callout}

In investing, return and risk must be analyzed together.

Return tells you how much a portfolio may grow.

Risk tells you how uncertain, unstable, or fragile that growth may be.

Two portfolios can have similar average returns but very different downside behavior.

\begin{callout}
\textbf{Example.}

\begin{itemize}
\item Portfolio A may grow more steadily with smaller drawdowns.
\item Portfolio B may have higher upside but much deeper losses in weak scenarios.
\end{itemize}
\end{callout}

\textbf{This simulator helps visualize both.}

\begin{itemize}
\item upside potential
\item downside tail risk
\item uncertainty of final outcomes
\end{itemize}

\section{Why Monte Carlo Simulation Is Useful}
\begin{callout}
\textbf{Level.} Beginner
\end{callout}

A single forecast is often misleading because markets are uncertain.

Monte Carlo simulation addresses this by generating many possible future paths instead of only one path.

\textbf{This helps answer questions such as.}

\begin{itemize}
\item What is the range of plausible outcomes?
\item How often do poor outcomes occur?
\item How severe are bad outcomes?
\item How uncertain is final wealth?
\end{itemize}

This is especially useful in portfolio analysis because investors do not only care about the average case.

They also care about tail risk, dispersion, and robustness.

\section{How the Simulation Works in Plain English}
\begin{callout}
\textbf{Level.} Beginner
\end{callout}

\textbf{At a high level, the app works like this.}

\begin{enumerate}
\item It downloads historical price data.
\item It converts prices into returns.
\item It estimates statistical inputs such as expected return, volatility, and covariance.
\item It generates many random future asset paths.
\item It combines those simulated asset paths using your chosen portfolio weights.
\item It summarizes the simulated outcomes using charts and metrics.
\end{enumerate}

\begin{callout}
\textbf{Important.}

This process is model-based.

It is useful for scenario exploration, learning, and structured risk analysis.

It is not a guarantee of future performance.
\end{callout}

\chapter{Reading the Outputs}

\section{How to Read the Simulated Portfolio Paths Chart}
\begin{callout}
\textbf{Level.} Beginner
\end{callout}

This chart shows many possible portfolio trajectories over time.

\textbf{How to interpret it.}

\begin{itemize}
\item each line represents one simulated future
\item a narrow cloud suggests more stable outcomes
\item a wide cloud suggests more uncertainty
\item downward paths reveal possible drawdown scenarios
\item upward paths illustrate upside potential
\end{itemize}

\textbf{What to look for.}

\begin{itemize}
\item the overall spread of paths
\item whether severe downside scenarios appear
\item whether growth appears steady or highly unstable
\item whether uncertainty widens quickly as the horizon increases
\end{itemize}

\section{How to Read the Return Distribution}
\begin{callout}
\textbf{Level.} Beginner
\end{callout}

The return distribution summarizes simulated portfolio returns.

\textbf{It helps you see.}

\begin{itemize}
\item the center of likely outcomes
\item the spread of outcomes
\item the downside tail
\item the upside tail
\end{itemize}

\textbf{Interpretation.}

\begin{itemize}
\item a distribution shifted to the right suggests stronger expected performance
\item a wider distribution suggests more uncertainty
\item a heavy left tail suggests larger downside risk
\item an asymmetric distribution may indicate uneven upside versus downside behavior
\end{itemize}

\section{How to Read the Terminal Value Distribution}
\begin{callout}
\textbf{Level.} Beginner
\end{callout}

The terminal value distribution focuses on ending wealth.

\textbf{This is important because many investors care more about.}

\begin{itemize}
\item what their portfolio may be worth at the end
\item how bad poor endings may be
\item how wide the gap is between weak and strong scenarios
\end{itemize}

\textbf{Interpretation.}

\begin{itemize}
\item the middle mass shows typical outcomes
\item the left tail shows poor endings
\item the right tail shows strong upside scenarios
\item a wide distribution means final wealth is highly uncertain
\end{itemize}

\section{Understanding Expected Return and Volatility}
\begin{callout}
\textbf{Level.} Intermediate depth
\end{callout}

Two of the most fundamental portfolio statistics are expected return and volatility.

\textbf{Expected return.}

\begin{equation}
\mathbb{E}[R_p] = w^\top \mu
\end{equation}

\textbf{Volatility.}

\begin{equation}
\sigma_p = \sqrt{w^\top \Sigma w}
\end{equation}

\textbf{where.}

\begin{itemize}
\item $w$ is the portfolio weight vector
\item $\mu$ is the expected return vector
\item $\Sigma$ is the covariance matrix
\end{itemize}

\textbf{Interpretation.}

\begin{itemize}
\item higher expected return suggests stronger average modeled growth
\item higher volatility means outcomes fluctuate more and are less predictable
\end{itemize}

A portfolio should not be judged by return alone.

The return-risk trade-off is central.

\section{Understanding VaR and CVaR}
\begin{callout}
\textbf{Level.} Intermediate depth
\end{callout}

Value at Risk (VaR) measures a downside threshold.

\textbf{For return $R$ and confidence level $\alpha$.}

\begin{equation}
\mathrm{VaR}_{\alpha}(R) = Q_{1-\alpha}(R)
\end{equation}

\begin{callout}
\textbf{Example.}

If 95\% VaR is -3\%, then only about 5\% of simulated outcomes are worse than -3\%.
\end{callout}

\textbf{Conditional Value at Risk (CVaR) goes further.}

\begin{equation}
\mathrm{CVaR}_{\alpha}(R) = \mathbb{E}[R \mid R \le \mathrm{VaR}_{\alpha}(R)]
\end{equation}

\textbf{Interpretation.}

\begin{itemize}
\item VaR gives the cutoff for bad outcomes
\item CVaR gives the average severity once you are already in that bad tail
\end{itemize}

CVaR is often more informative because it captures tail depth, not only the threshold.

\chapter{Methodology}

\section{Portfolio Construction Mathematics}
\begin{callout}
\textbf{Level.} Technical
\end{callout}

\textbf{Let the portfolio weights be.}

\begin{equation}
w = (w_1, w_2, \dots, w_n)
\end{equation}

\textbf{For a fully invested portfolio.}

\begin{equation}
\sum_{i=1}^{n} w_i = 1
\end{equation}

\textbf{Portfolio return.}

\begin{equation}
R_p = \sum_{i=1}^{n} w_i R_i
\end{equation}

\textbf{Expected portfolio return.}

\begin{equation}
\mathbb{E}[R_p] = w^\top \mu
\end{equation}

\textbf{Portfolio variance.}

\begin{equation}
\mathrm{Var}(R_p) = w^\top \Sigma w
\end{equation}

These formulas form the mathematical foundation of portfolio analysis in the app.

\textbf{Key insight.}

Even if two assets are risky individually, combining them can reduce overall portfolio risk if their movements are not perfectly aligned.

\section{Geometric Brownian Motion (GBM) Foundation}
\begin{callout}
\textbf{Level.} Technical
\end{callout}

\textbf{The simulator models asset prices using Geometric Brownian Motion (GBM).}

\begin{equation}
dS_t = \mu S_t \, dt + \sigma S_t \, dW_t
\end{equation}

\textbf{where.}

\begin{itemize}
\item $S_t$ is asset price at time $t$
\item $\mu$ is drift
\item $\sigma$ is volatility
\item $W_t$ is Brownian motion
\end{itemize}

\textbf{A common discrete-time simulation form is.}

\begin{equation}
S_{t+\Delta t} = S_t \exp\left[
\left(\mu - \frac{1}{2}\sigma^2\right)\Delta t
+ \sigma \sqrt{\Delta t} Z
\right]
\end{equation}

\textbf{where.}

\begin{itemize}
\item $\Delta t$ is the time step
\item $Z \sim N(0,1)$
\end{itemize}

\textbf{Why GBM is commonly used.}

\begin{itemize}
\item it keeps prices positive
\item it is mathematically tractable
\item it scales well for simulation
\item it is a standard baseline model in finance
\end{itemize}

\textbf{Why it is imperfect.}

\begin{itemize}
\item it assumes constant parameters
\item it may understate extreme events
\item it does not capture all real market dynamics
\end{itemize}

\section{Parameter Estimation}
\begin{callout}
\textbf{Level.} Technical
\end{callout}

The simulator estimates model inputs from historical return data.

\textbf{Typical estimated quantities.}

\begin{itemize}
\item expected return vector $\mu$
\item volatility
\item covariance matrix $\Sigma$
\end{itemize}

\textbf{These estimates drive.}

\begin{itemize}
\item simulation
\item portfolio risk metrics
\item optimization
\item projection outputs
\end{itemize}

\textbf{Important caution.}

Historical estimates are noisy.

Future market behavior may differ from the past.

\textbf{This means the app should be interpreted as.}

\begin{itemize}
\item a structured approximation
\item a decision-support tool
\item a learning framework
\end{itemize}

not as a certainty engine.

\section{Monte Carlo Engine Logic}
\begin{callout}
\textbf{Level.} Technical
\end{callout}

Monte Carlo simulation generates many random future scenarios.

\textbf{For each simulation path, the engine.}

\begin{enumerate}
\item draws random shocks
\item propagates asset prices forward using the model
\item computes portfolio values using the selected weights
\item records path-level and terminal outcomes
\end{enumerate}

If there are $M$ simulated paths, the output is an empirical distribution across those $M$ scenarios.

\textbf{This allows the app to estimate.}

\begin{itemize}
\item expected return
\item volatility
\item quantiles
\item VaR and CVaR
\item terminal wealth percentiles
\item projection outcomes under repeated contributions
\end{itemize}

The power of Monte Carlo lies in studying the distribution of outcomes, not only a single expected value.

\section{Portfolio Optimization}
\begin{callout}
\textbf{Level.} Intermediate depth
\end{callout}

Portfolio optimization studies how different allocations change the return-risk trade-off.

\textbf{The app highlights common candidate portfolios such as.}

\begin{itemize}
\item current portfolio
\item maximum Sharpe portfolio
\item minimum volatility portfolio
\end{itemize}

\textbf{Sharpe ratio.}

\begin{equation}
\mathrm{Sharpe} = \frac{\mathbb{E}[R_p] - R_f}{\sigma_p}
\end{equation}

\textbf{where.}

\begin{itemize}
\item $\mathbb{E}[R_p]$ is expected portfolio return
\item $R_f$ is the risk-free rate
\item $\sigma_p$ is portfolio volatility
\end{itemize}

\textbf{Interpretation.}

\begin{itemize}
\item higher Sharpe suggests stronger risk-adjusted return
\item minimum volatility focuses on stability
\item your current portfolio may or may not be efficient relative to alternatives
\end{itemize}

Optimization does not reveal one universally correct portfolio.

It reveals how assumptions and allocations shape the opportunity set.

\section{How to Read the Efficient Frontier}
\begin{callout}
\textbf{Level.} Intermediate depth
\end{callout}

\textbf{The efficient frontier is the set of portfolios that offer.}

\begin{itemize}
\item the highest expected return for a given volatility, or
\item the lowest volatility for a given expected return
\end{itemize}

\textbf{The app approximates this by generating many candidate portfolios and plotting.}

\begin{itemize}
\item x-axis: volatility
\item y-axis: expected return
\end{itemize}

\textbf{Highlighted portfolios often include.}

\begin{itemize}
\item Current Portfolio
\item Max Sharpe Portfolio
\item Min Volatility Portfolio
\end{itemize}

\textbf{Interpretation.}

\begin{itemize}
\item points further left are lower risk
\item points higher up are higher return
\item points below attractive alternatives may be inefficient under the model assumptions
\end{itemize}

The frontier helps users compare design choices rather than focusing only on one allocation.

\chapter{Validation and Extensions}

\section{Historical Backtesting}
\begin{callout}
\textbf{Level.} Intermediate depth
\end{callout}

Backtesting applies the chosen portfolio weights to realized historical returns.

\textbf{This helps answer.}

\begin{itemize}
\item how the portfolio would have behaved historically
\item whether it outperformed or underperformed a benchmark
\item whether historical performance was smooth or volatile
\end{itemize}

\textbf{A common cumulative growth representation is.}

\begin{equation}
V_t = \prod_{s=1}^{t} (1 + R_s)
\end{equation}

where $R_s$ is the portfolio return in period $s$.

Backtesting is useful because it adds historical context.

But it must not be confused with future certainty.

A strategy that looked strong historically may still perform poorly going forward.

\section{Benchmark Comparison}
\begin{callout}
\textbf{Level.} Beginner
\end{callout}

A benchmark provides context for performance.

A portfolio gaining 8\% may sound good in isolation.

But if a simple benchmark gained 15\% with similar or lower risk, the result looks less impressive.

\textbf{Common reasons to compare against a benchmark.}

\begin{itemize}
\item evaluate active choices fairly
\item determine whether complexity added value
\item understand if performance is mostly just market exposure
\end{itemize}

A broad-market ETF such as SPY is often used as a reference benchmark.

\section{CAPM / Market Factor Analysis}
\begin{callout}
\textbf{Level.} Technical
\end{callout}

\textbf{A simple market factor model is.}

\begin{equation}
R_p - R_f = \alpha + \beta (R_m - R_f) + \varepsilon
\end{equation}

\textbf{where.}

\begin{itemize}
\item $R_p$ is portfolio return
\item $R_m$ is market return
\item $R_f$ is risk-free rate
\item $\alpha$ is alpha
\item $\beta$ is beta
\item $\varepsilon$ is the residual component
\end{itemize}

\textbf{Interpretation.}

\begin{itemize}
\item $\beta$ measures market sensitivity
\item $\alpha$ measures return not explained by market exposure alone
\item $R^2$ measures how much of return variation is explained by the model
\end{itemize}

\textbf{This helps distinguish between.}

\begin{itemize}
\item market-driven behavior
\item portfolio-specific behavior
\end{itemize}

\section{How to Read the CAPM Scatter Plot}
\begin{callout}
\textbf{Level.} Intermediate depth
\end{callout}

Each point in the CAPM scatter plot represents one observation period.

\textbf{Axes.}

\begin{itemize}
\item x-axis: market return
\item y-axis: portfolio return
\end{itemize}

The fitted line summarizes the average relationship.

\textbf{Interpretation.}

\begin{itemize}
\item steeper slope means higher beta
\item tighter clustering around the line means higher R²
\item wider scatter means more idiosyncratic behavior
\item a positive intercept suggests positive alpha under the fitted model
\end{itemize}

This chart helps users understand whether the portfolio behaves mostly like the market or has distinct return dynamics.

\section{Stress Testing}
\begin{callout}
\textbf{Level.} Intermediate depth
\end{callout}

Stress testing asks how the portfolio behaves under adverse assumptions.

\textbf{Examples.}

\begin{itemize}
\item lower expected returns
\item higher volatility
\item market shock scenarios
\item crash-style drawdowns
\item worse-than-normal environments
\end{itemize}

\textbf{Why it matters.}

A portfolio can appear strong in normal conditions but fail badly under stress.

\textbf{Stress testing helps reveal.}

\begin{itemize}
\item fragility
\item downside sensitivity
\item robustness
\item concentration risk
\item dependence on favorable assumptions
\end{itemize}

\section{Investment Outcome Projections}
\begin{callout}
\textbf{Level.} Beginner
\end{callout}

This section translates portfolio return assumptions into practical long-term investing outcomes.

\textbf{Typical inputs.}

\begin{itemize}
\item initial investment
\item monthly contribution
\item time horizon
\end{itemize}

\textbf{Typical outputs.}

\begin{itemize}
\item median final value
\item mean final value
\item pessimistic outcomes
\item optimistic outcomes
\item range of plausible final wealth
\end{itemize}

This makes the app more practical because investors often care about future wealth paths, not just abstract return distributions.

\section{How to Interpret Projection Results}
\begin{callout}
\textbf{Level.} Beginner
\end{callout}

\textbf{Median Final Value.}

The middle outcome. Half of simulations end above it and half below it.

\textbf{Mean Final Value.}

The average outcome. This can be pulled upward by a relatively small number of strong scenarios.

\textbf{5th Percentile.}

A pessimistic threshold representing weak outcomes.

\textbf{95th Percentile.}

An optimistic threshold representing strong outcomes.

\textbf{Best practice.}

\begin{itemize}
\item use the median as a practical central case
\item use low percentiles for risk planning
\item use high percentiles to understand upside potential
\end{itemize}

\chapter{Cautions and Practice}

\section{Common Portfolio Analysis Mistakes}
\begin{callout}
\textbf{Level.} Beginner
\end{callout}

\textbf{Common mistakes include.}

\begin{itemize}
\item focusing only on expected return
\item ignoring downside tail risk
\item treating historical averages as certainty
\item overtrusting optimization output
\item ignoring benchmark comparison
\item forgetting concentration risk
\item assuming a smooth chart means low real-world risk
\item confusing model output with financial advice
\end{itemize}

\textbf{A better habit is to combine.}

\begin{itemize}
\item simulation
\item risk metrics
\item historical context
\item stress testing
\item judgment
\end{itemize}

\section{Model Assumptions and Limitations}
\begin{callout}
\textbf{Level.} Technical
\end{callout}

\textbf{Important limitations of this framework include.}

\begin{itemize}
\item GBM assumes constant drift and volatility
\item correlations may change over time
\item historical estimates can be unstable
\item regime shifts are difficult to capture
\item extreme events may be more severe than modeled
\item transaction costs may be ignored
\item taxes may be ignored
\item rebalancing rules may not be fully modeled
\item future market structure may differ from the estimation window
\end{itemize}

\textbf{Therefore, the simulator should be used as.}

\begin{itemize}
\item a learning tool
\item a scenario engine
\item a structured approximation
\end{itemize}

It should not be interpreted as certainty or personal financial advice.

\section{FAQ}
\begin{callout}
\textbf{Level.} Beginner
\end{callout}

\begin{callout}
\textbf{Q.} Does this app predict the future exactly?

\textbf{A.} No. It generates many possible futures under model assumptions.
\end{callout}

\begin{callout}
\textbf{Q.} Why do results change from run to run?

\textbf{A.} Because random simulation is part of Monte Carlo analysis.
\end{callout}

\begin{callout}
\textbf{Q.} Is higher expected return always better?

\textbf{A.} No. Risk, dispersion, and tail behavior also matter.
\end{callout}

\begin{callout}
\textbf{Q.} Why compare with SPY?

\textbf{A.} It provides market context.
\end{callout}

\begin{callout}
\textbf{Q.} Should I trust the max-Sharpe portfolio automatically?

\textbf{A.} No. Optimization depends heavily on estimated inputs and assumptions.
\end{callout}

\begin{callout}
\textbf{Q.} Why are VaR and CVaR useful?

\textbf{A.} They summarize downside threshold and tail severity.
\end{callout}

\begin{callout}
\textbf{Q.} Does strong backtest performance guarantee future success?

\textbf{A.} No. Historical evidence is informative, but never a guarantee.
\end{callout}

\begin{callout}
\textbf{Q.} Why can terminal wealth vary so much?

\textbf{A.} Because compounding magnifies differences over time, especially under volatility.
\end{callout}

\end{document}
